% -----------------------------------------------------------------------------------------------------
\section{超出领头幂次的 TMD 因子化}\label{sec:front-nlp-tmd}

Drell--Yan、SIDIS 与 $e^+e^-$ 的领头幂次 TMD 因子化是全阶论证（\cref{thm:tmd-factorization}，[A]）。在 $q_T/Q$ 的次领头幂次，修正分三类：运动学的（领头项的精确运动学与参考系依赖）、内禀的／扭度 3 的（夸克—胶子—夸克关联，“TMD 扭度”）与动力学的（新的硬函数）。下面所有结果都是微扰的，并且或明或暗地假定 Glauber 交换在 NLP 不破坏因子化。

\begin{conjecture}[次领头幂次 TMD 过程中 Glauber 交换的相消 \tagC]\label{conj:front-glauber-nlp}
在 Drell--Yan 与 SIDIS 中 $q_T/Q$ 的次领头幂次，Glauber 交换相消到足以使\cref{prop:front-tmd-nlp-operator,prop:front-sidis-nlp}的 NLP 因子化公式对强子截面成立的程度（领头幂次相消\cref{lem:tmd-glauber}的类比）。
\uses{def:tmd-kin, thm:tmd-factorization}
\end{conjecture}

\paragraph{地位与缺口。}[C]：处处被假定——\cite{Ebert:2021jhy}明言“working under the assumption that leading power Glauber interactions do not spoil factorization at this order”，其他工作则隐含地假定；在 NLP 没有全阶论证。陈述中对后面两个命题的引用只是文字上的指向，不是依赖。

\begin{proposition}[TMD 算符展开与运动学幂次修正 \tagA]\label{prop:front-tmd-nlp-operator}
以下陈述针对 TMD 算符展开本身（只含共线与软自由度），它在强子截面中的有效性另需\cref{conj:front-glauber-nlp}。到 NLO 与 NLP：(a) TMD 因子化在算符层次成立，算符按 TMD 扭度排序（分开运动学修正与真正的幂次修正），过程依赖经由背景场的边界条件进入；对任何有两个被探测态的过程，强子张量的有效算符构造到 NLP（$\sim q_T/Q$）与 NLO\cite{Vladimirov:2021hdn}；(b) 扭度 3 TMD 的定义、演化方程（大 $N_c$ 下自治）、小 $b$ 渐近、奇异结构与“物理的”扭度 3 TMD\cite{Rodini:2022wki}（连同勘误）；(c) 每个真正的项都生成一串与之有相同非微扰内容的运动学幂次修正，领头项所带的一串可以求和，恢复在固定幂次上被破坏的电磁流守恒与参考系无关性，得到类似自由无质量夸克强子张量的因子化形式，并在 NLO 核对\cite{Vladimirov:2023aot}；SIDIS 版本见\cite{Piloneta:2025jjb}，准 TMD 版本（联系\cref{thm:lamet-quasi-tmd}）见\cite{BrisCuerpo:2026pma}（以上均据各文摘要）。（本命题只涉及算符展开；强子层次另需\cref{asm:intro-bridge}与上述猜想。）
% errtype: E-S
\uses{thm:tmd-factorization, def:tmd-collins}
\end{proposition}

\begin{proofsketch}
在背景场方法中把强子张量算符按 TMD 扭度展开，逐阶（NLO）匹配系数，并对运动学修正的级数求和。完整论证见所引文献。
\end{proofsketch}

\paragraph{地位与缺口。}[A]（NLO、NLP；只读到摘要层次）。首句中对\cref{conj:front-glauber-nlp}的引用不是依赖：陈述限于不含 Glauber 的算符展开。

\begin{proposition}[SIDIS 方位角不对称的次领头幂次 SCET 因子化 \tagA]\label{prop:front-sidis-nlp}
在不含 Glauber 模式的 SCET 中（Glauber 相消见\cref{conj:front-glauber-nlp}），SIDIS 方位角不对称中幂次压低的硬散射效应在 NLP 完全因子化：出现依赖两个纵向动量分数的 TMD（其中一个只在树图以上出现）；在有效理论内只出现领头幂次软函数，并被吸收进 TMD；只需一个新的硬系数（\cite{Ebert:2021jhy}，连同勘误；据该文摘要）。（本命题只涉及有效理论中的因子化结构。）
% errtype: E-S
\uses{thm:tmd-factorization, def:tmd-soft, def:scet-modes, asm:front-scet-modes}
\end{proposition}

\begin{proofsketch}
在 SCET 中构造 NLP 的流与拉氏量插入，证明软函数只以领头幂次形式出现并可吸收，再与 QCD 匹配。完整论证见\cite{Ebert:2021jhy}。
\end{proofsketch}

\paragraph{地位与缺口。}[A]，限于所述条件；强子陈述需要\cref{conj:front-glauber-nlp}（[C]）。用到\cref{asm:front-scet-modes}的 (SM1)，(SM2) 由陈述中的限制代替。

\begin{proposition}[单圈次领头幂次系数函数与扭度 3 TMD 演化 \tagP]\label{prop:front-nlp-tmd-oneloop}
在单圈阶（均据各文摘要）：(a) SIDIS 全部结构函数在 NLP 的单圈系数函数与演化核、扭度 3 TMD 碎裂函数与领头阶核——作者指出与早先推导有“several differences”，源于夸克—胶子干涉项的运动学\cite{Rodini:2023plb}；(b) Drell--Yan 与 SIDIS 的运动学、内禀与动力学 NLP 部分，单圈的重整化群相容性，混合领头与次领头分布的反常维度矩阵，与领头幂次不同的硬函数和软函数，动力学 NLP 分布的快度反常维度等于领头幂次的\cite{Gamberg:2022lju}（预印本）；(c) 在不显含软、共线模式的有效理论中，NLP 的 Drell--Yan：幂次压低算符 $T_{(i,j)}$ 给出普通 PDF 的 $\order{q_T^2/q^2}$ 系数，$q^2\gg q_T^2\gg\Lambda^2$ 时的 NLP 因子化，快度发散的相消与重叠减除\cite{Inglis-Whalen:2021bea}；(d) Wilson 线方法：TMD 因子化的幂次修正\cite{Balitsky:2017flc,Balitsky:2017gis}，以及大 $N_c$ 下保持电磁规范不变的非极化 Drell--Yan 张量的 $1/Q^2$ 修正\cite{Balitsky:2024ozy}。（本命题只涉及部分子系数函数与演化核。）
\uses{prop:front-tmd-nlp-operator}
\end{proposition}

\paragraph{地位与缺口。}[P]（单圈）；各组的结果尚不完全一致，我们照各文的陈述报告，不作裁决。

\begin{remark}[次领头幂次的幂律快度发散]\label{rem:front-power-rapidity}
据该文摘要\cite{Ebert:2018gsn}，在次领头幂次快度发散可以是幂律型的，并产生 PDF 的导数；纯快度正规子与类 $\MSbar$ 方案保持幂次展开的齐次性；色单态产生的全部 $\order{\alpha_s}$ $q_T^2/Q^2$ 修正已经给出。超出单圈与非微扰的 NLP TMD 因子化仍是开放的。
\end{remark}
